% ECONOMICS_PROBLEM: {"id": "D82-92", "title": "Multidimensional screening", "jel_code": "D82", "source_id": "OP-0006"}
% !TeX program = xelatex
\documentclass[11pt,a4paper]{article}
\usepackage[margin=23mm]{geometry}
\usepackage{fontspec}
\setmainfont{Latin Modern Roman}
\usepackage{amsmath,amssymb,mathtools}
\usepackage{xcolor,enumitem,booktabs,longtable,array,xurl,hyperref}
\definecolor{muted}{HTML}{53636C}
\hypersetup{colorlinks=true,urlcolor=blue,linkcolor=blue}
\setlength{\parindent}{0pt}
\setlength{\parskip}{5pt}
\newcommand{\R}{\mathbb R}
\newcommand{\E}{\mathbb E}
\newcommand{\N}{\mathbb N}
\newcommand{\ind}{\mathbf 1}
\newcommand{\Var}{\operatorname{Var}}
\newcommand{\Cov}{\operatorname{Cov}}
\newcommand{\supp}{\operatorname{supp}}
\DeclareMathOperator*{\argmin}{arg\,min}
\DeclareMathOperator*{\argmax}{arg\,max}
\newcommand{\entrymeta}[3]{{\sffamily\small\color{muted}\textbf{\texttt{#1}}\quad|\quad #2\par #3\par}}
\newcommand{\Problemref}[1]{\texttt{#1}}
\newcommand{\JELref}[2]{\texttt{#2} (JEL \texttt{#1})}
\begin{document}
\section*{Multidimensional screening}
\textbf{Problem ID:} \texttt{D82-92}\quad \textbf{JEL:} \texttt{D82}\par
% BEGIN ECONOMICS STATEMENT
\label{problem:OP-0006}
\label{atlas:prob:M6}

\noindent\textbf{Original atlas ID:} M6. \textbf{Formulation:} editorial specification of a research family.

\subsection*{Definitions and assumptions}

Let the consumer type space $\Theta\subset\mathbb R^d$, $d\ge2$, be compact and convex, with a specified continuously positive density $f$ relative to Lebesgue measure. Let $Q\subset\mathbb R_+^d$ be a compact convex feasible allocation set, let $C:Q\to\mathbb R$ be a continuous convex production cost, and normalize outside-option utility to zero. A direct mechanism consists of measurable $q:\Theta\to Q$ and payment $t:\Theta\to\mathbb R$ with integrable revenue. Utility is quasilinear: $U(\theta)=\theta\cdot q(\theta)-t(\theta)$. Incentive compatibility requires $U(\theta)\ge\theta\cdot q(\hat\theta)-t(\hat\theta)$ for every $\theta,\hat\theta\in\Theta$; individual rationality requires $U(\theta)\ge0$ for every type. These are global, not merely local, constraints.

\subsection*{Formal research question}

For a declared distribution-and-cost class $\mathcal C$, characterize the maximizers, or certified approximate maximizers, of
\[
\sup_{(q,t)\;\mathrm{IC,IR}}\int_\Theta[t(\theta)-C(q(\theta))]f(\theta)\,\mathrm d\theta.
\]
Equivalently, impose convexity of $U$, $q(\theta)\in\partial U(\theta)\cap Q$ for every type, and $t(\theta)=\theta\cdot q(\theta)-U(\theta)$. The subgradient here is defined by $U(\hat\theta)\ge U(\theta)+q(\theta)\cdot(\hat\theta-\theta)$ for all $\hat\theta\in\Theta$. Determine conditions for exclusion and bunching, and develop computation that proves global feasibility and an objective gap no larger than $\varepsilon>0$. State approximation rates and dependence on dimension and primitive regularity, rather than reporting only agreement on a finite type grid.

\subsection*{Interpretation and scope}

In this linear-type model the convex-potential conditions are equivalent to the full incentive inequalities. The equivalence does not automatically extend to arbitrary nonlinear preferences, random outside options, or multiple competing principals. Bunching refers to loss of type distinctions over a positive-measure region of the type space. Where $q$ is differentiable, a precise benchmark is a measurable region $B\subseteq\Theta$ of positive Lebesgue measure with $\operatorname{rank}Dq(\theta)<d$ almost everywhere on $B$. An exact common allocation for a positive-measure set is a stronger pooling special case, not the general definition. Existence, uniqueness, and geometry must be established for the particular class $\mathcal C$ and cannot be inferred from a discretized numerical optimum. Continuous densities, compact quantity bounds, and a normalization such as $\min U=0$ help control the problem, but a complete computational guarantee still needs a valid discretization or dual bound.

\subsection*{Source audit and status}

Rochet--Choné (1998), Armstrong (1996), and Rochet--Stole (2003) are verified; the latter is chapter 5 of the Eighth World Congress volume, printed in 2003, not its 2010 online-release year. Rochet--Choné already supplies existence and characterization for a particular model. Original link [12] is a person profile rather than a research article and does not substantiate the open-status claim; it is replaced by actual publications. Rochet's 2024 primary review [S4] confirms substantial subsequent development. The specified optimization and certification agenda is editorial; neither multidimensional screening as a whole nor its foundational benchmark should be labeled wholly unsolved.

\subsection*{References}

\begin{enumerate}

\item[S1] Jean-Charles Rochet and Philippe Choné (1998). \emph{Ironing, Sweeping, and Multidimensional Screening}. Econometrica 66(4), 783--826. \url{https://researchportal.ip-paris.fr/en/publications/ironing-sweeping-and-multidimensional-screening/}. \textit{Locator:} abstract; multiproduct monopoly characterization. \textit{Establishes:} Existence and characterization for a class of multidimensional screening problems. \textit{Does not establish:} an unsolved status for its already characterized benchmark.

\item[S2] Mark Armstrong (1996). \emph{Multiproduct Nonlinear Pricing}. Econometrica 64(1), 51--75. \url{https://www.jstor.org/stable/2171924}. \textit{Locator:} article; multidimensional pricing analysis. \textit{Establishes:} A tractable multidimensional nonlinear-pricing approach. \textit{Does not establish:} a complete characterization for arbitrary joint type distributions.

\item[S3] Jean-Charles Rochet and Lars A. Stole (2003). \emph{The Economics of Multidimensional Screening}. Advances in Economics and Econometrics: Theory and Applications, Eighth World Congress, Volume 1, chapter 5, 150--197; Cambridge University Press. \url{https://doi.org/10.1017/CBO9780511610240.006}. \textit{Locator:} chapter 5; introduction. \textit{Establishes:} A survey of multidimensional screening and economic applications. \textit{Does not establish:} current certification that every multidimensional screening problem is open.

\item[S4] Jean-Charles Rochet (2024). \emph{Multidimensional Screening after 37 Years}. Journal of Mathematical Economics 113, article 103010. \url{https://www.sciencedirect.com/science/article/pii/S0304406824000715}. \textit{Locator:} abstract; sections on rationalizability and monopoly pricing. \textit{Establishes:} A recent account of implementability and multidimensional screening developments. \textit{Does not establish:} that no progress occurred after the 1998 benchmark.

\end{enumerate}

\subsection*{Common conventions: Source mathematical conventions}

Unless an entry states otherwise, agent and firm sets are finite. State and action spaces are measurable, strategies and outcome maps are measurable, probability laws are specified, and expectations exist for the targets under discussion. Infinite-horizon discount factors lie in $(0,1)$ and discounted objectives are finite. Norms, tolerances and complexity input sizes must be specified for computational comparisons. Constants or primitives that appear locally are inputs, not empirical facts asserted by this catalogue.

\textbf{Identification.} A statistical model $m\in\mathcal M$ implies an observable law $P_m$ and target $\theta(m)$. At law $P$, its identified set is
\[
\Theta(P;\mathcal M)=\{\theta(m):m\in\mathcal M,\ P_m=P\}.
\]
Point identification holds when this set is a singleton, or equivalently when $P_{m_1}=P_{m_2}$ implies $\theta(m_1)=\theta(m_2)$ for all compatible models. Sharp bounds mean bounds attained, or approached when the boundary is not attained, by compatible models. Writing this set is a definition, not a solution: the substantive task is to establish informative restrictions and derive or estimate the set. An unrestricted-confounding model may give only trivial bounds.

\textbf{Inference.} Identification, estimability and valid uncertainty quantification are separate questions. Fix $\alpha\in(0,1)$. A confidence procedure $C_n$ over a declared law class $\mathcal P$ has asymptotic uniform coverage at level $1-\alpha$ when
\[
\liminf_{n\to\infty}\inf_{P\in\mathcal P}\Pr_P\{\theta(P)\in C_n\}\geq1-\alpha.
\]
For set-identified targets the coverage event must be defined separately, for example coverage of the full identified set. If an entry uses identification-indexed models instead of uniquely indexed laws, the same coverage convention is applied over those models. Pointwise limits do not establish uniform validity, and asymptotic validity does not guarantee finite-sample precision. Sample sizes, dependence structures and weak-identification sequences are part of each application.

\textbf{Counterfactuals.} $Y_i(a)$ is a potential outcome under a specified intervention, including its timing, implementation and versions. It is not $Y_i$ conditional on voluntarily choosing $a$. A full intervention vector permits interference; shorthand scalar treatment notation does not silently impose no interference. Assignment, selection, attrition, anticipation and measurement must be specified. A claim about an instrument's local effect is not automatically a population policy effect.

\textbf{Economic equilibrium.} A counterfactual model includes best responses, constraints, feasibility, market clearing, state transitions and expectations consistent with the stated information. When several equilibria exist, either a declared selector is an input or the target is set-valued. Comparisons across policy regimes must use the stated selection convention. Reduced-form relationships are not presumed invariant after an equilibrium-changing intervention.

\textbf{Optimization and welfare.} Optimization is over the stated feasible set. If attainment is not established, $\sup$ or $\inf$ and $\varepsilon$-optimal choices replace an asserted maximizer or minimizer. For example, with value $V=\sup_{a\in\mathcal A}W(a)<\infty$, an $\varepsilon$-optimal set is $\{a\in\mathcal A:W(a)\geq V-\varepsilon\}$ for $\varepsilon>0$. Benefits, costs, risk and health or environmental outcomes can be aggregated only using declared units and conversion weights. Interpersonal utility weights, the treatment of fiscal transfers and foreign profits, discounting, and the behavioral welfare criterion are explicit inputs. A comparative-static derivative additionally requires existence and appropriate differentiability of the selected counterfactual.

\textbf{Sources.} Local labels [S1], [S2], and so forth restart in every question. Locators identify the relevant abstract, model, section or result at the level actually checked; a bibliographic or abstract-level verification is not represented as inspection of an entire proof. Working papers are distinguished from later publications and changing versions. Source summaries are paraphrased. All URL checks and editorial formulations refer to the audit date stated above.
% END ECONOMICS STATEMENT
\end{document}
